% ============================================================================
% sec1-3.tex — satellite of main_submission.tex (§1 Introduction, §2
% Preliminaries + campaign notation, §3 Fan-type shell). Split verbatim from
% the batch-1 single file by deep_P14 batch 2; content unchanged.
% Ownership: this file carries the batch-1 material; further edits belong to
% the Abstract/§1 rewrite batch (P-ARCH S8), not to campaign-chapter batches.
% Migration anchors are inline (main.tex line refs = ledger state 8b5323c5…).
% ============================================================================
% ============================================================================
% §1 INTRODUCTION (rewritten per P-ARCH §2.2: literature survey retained,
% honesty note retained verbatim, Contributions rearranged 4 -> 5 blocks
% (machine-campaign package and conditional full chain added), three-tier
% declaration as a formal paragraph, six no-claims wording in the body).
% Migration anchors: source §1 main.tex :93-261 (survey :101-121, honesty
% note :123-128, semi-clique paragraph :130-135, contributions :137-228,
% methodology :230-237, adjudication paragraph :239-253, open problems
% :255-260); six no-claims verbatim register main.tex :5378-5396.
% ============================================================================
\section{Introduction}\label{sec:intro}

Let $G$ be a finite connected simple graph on $n$ vertices. A \emph{path decomposition} of $G$
is a partition of $E(G)$ into paths, and $\pn(G)$ denotes the minimum number of paths in such a
decomposition. Gallai conjectured that $\pn(G)\le \ceil{n/2}$ for every connected $G$; the
conjecture was communicated by Erd\H{o}s and is tracked as problem \#583 in \cite{EP583}. It
remains open, with $\floor{2n/3}$ the best general bound, due to Dean and Kouider \cite{DK00}.

The case analysis of the last six decades has been organized, to a remarkable degree, by the
structure of the subgraph $\EV(G)$ induced by the even-degree vertices, together with bounds on
the degrees of $G$ itself. Lov\'asz \cite{Lo68} showed that $E(G)$ decomposes into at most
$\floor{n/2}$ paths and cycles, and settled the case $|\EV(G)|\le 1$; Pyber \cite{Pyb96} proved
that $\EV(G)$ a forest implies $\pn(G)\le \floor{n/2}$; Favaron and Kouider \cite{FK88} settled
Eulerian graphs of maximum degree~$4$; and Fan \cite{Fan05} proved the conjecture when every
block of $\EV(G)$ is triangle-free with maximum degree at most $3$. Two strengthenings followed:
Botler and Sambinelli \cite{BS21} proved it when $\EV(G)$ is a subgraph of a graph in their
class $\mathcal{G}$ (blocks of maximum degree at most $3$, at most one triangular block per
component; the strengthening being the class --- their main theorem is stated with Gallai's
ceiling $\ceil{n/2}$ rather than the floor of Theorem~\ref{thm:fan} below), and Fan, Hou and
Zhou \cite{FHZ23} when $\EV(G)$ contains no pair of triangles and
every block has maximum degree at most $3$. A parallel track bounds the degrees of $G$ itself:
maximum degree $\le 5$ \cite{BP19}, maximum degree $\le 6$ with an independence condition on
degree-$6$ vertices \cite{CFL}, treewidth $\le 3$ \cite{BCLS20}. Very recently the
class-based programme jumped forward: Zhang, Liu and Hong \cite{ZLH24} proved the conjecture,
with the sharper $\floor{n/2}$ bound, for all $3$-degenerate graphs except $K_3$ and $K_5^-$,
and Lu and Yi \cite{LY24} did the same for treewidth at most $4$ with exceptions $K_3$, $K_5^-$,
$K_5$; Chen and Wu \cite{CW25} settled it for block graphs. On the algorithmic side, the exact
computation of $\pn(G)$ has very recently been shown to be polynomial-time on subcubic graphs
\cite{MY26}.

\begin{quote}\small\itshape
Honesty note. All results of the present paper lie outside the classes of
\cite{ZLH24,LY24,CW25}, or are quantitative refinements of them; where our classes overlap
theirs we say so explicitly. In particular, our theorems in Section~\ref{sec:astar} require a
structural condition on $\EV(G)$ precisely where \cite{ZLH24,LY24} require none, and vice versa.
\end{quote}

Bonamy and Perrett \cite{BP19} also asked whether $\pn(G)\le \floor{n/2}$ for every connected
$G$ that is not an \emph{odd semi-clique} (a graph obtained from $K_{2k+1}$ by deleting at most
$k-1$ edges). The odd semi-cliques $K_3$ and $K_5^-$ are exactly the obstructions appearing in
\cite{ZLH24,LY24}, and this is forced by a two-line density arithmetic: a $3$-degenerate odd
semi-clique on $n=2k+1$ vertices has $2k^2+1\le 3n-6$ edges, whence $k\le 2$; cf.\ Question~1.1
of \cite{BP19}.

\paragraph{Contributions.}
\begin{enumerate}[label=(\roman*),leftmargin=2em]
\item \textbf{The exact path number of $K_{4,n}$ for odd $n$} (Section~\ref{sec:k4odd}).
Constantinou and Ellinas \cite{CE18} gave minimal path-decomposition algorithms for complete
bipartite graphs in three of four cases, leaving open the case of odd $n_1$ with
$1\le n_2\le n_1-1$; their concluding remark asks first for ``the exact characteristics'' of
this case. We close the $n_2=4$ column completely:
\[
\pn(K_{4,n})=(n+1)/2 \quad\text{for every odd } n\ge 3,
\]
which extends the even-case formula of \cite{CE18} into Case~4 and improves the best published
bound (a factor $k+2$ from \cite{Pyb96,Fan05,LY24}, where $n=2k+1$) by one. The proof has three
self-contained parts: a counting lower bound (each path uses at most $8$ edges), a structure
theorem pinning the length profile of every optimal decomposition, and an explicit closed-form
construction, supplemented by exhaustive small-case witnesses. For the remaining even columns
we prove two general tools: a structure theorem answering the ``exact characteristics'' request
of \cite{CE18} across the \emph{whole} even range (every optimal decomposition visits every
vertex of the small side, and a deficit ledger pins the length profile;
Theorem~\ref{thm:structgen}), and a \emph{ring-closure theorem}
(Theorem~\ref{thm:ringclosure}): in the column $n_2=2h$, one explicit base decomposition of
$K_{2k_0+1,2h}$ into $k_0+1$ paths extends, by an explicit construction, to every
$k\ge k_0+h+1$. This upgrades our conjecture (Conjecture~\ref{conj:27A}) into a
column-by-column decision procedure with $h+1$ base instances per column. As applications we
close the $n_2=6$ column completely,
\[
\pn(K_{n_1,6})=(n_1+1)/2 \quad\text{for every odd } n_1\ge 7,
\]
which settles the previously bracketed instance $K_{7,6}$ at $\pn=4$ and improves the best
published bound for that family by two; and in the columns $n_2=8$ and $n_2=10$ we prove six
further exact values, $\pn(K_{9,8})=5$, $\pn(K_{11,8})=6$, $\pn(K_{13,8})=7$,
$\pn(K_{11,10})=6$, $\pn(K_{13,10})=7$, $\pn(K_{15,10})=8$, together with the whole tails
$k\ge9$ resp.\ $k\ge11$ of those columns, leaving five named base instances open
(Section~\ref{sec:open}).
\item \textbf{Odd clique chains: refinements beyond block graphs} (Section~\ref{sec:chains}).
Chen and Wu \cite{CW25} proved Gallai's conjecture for all block graphs. For the subfamily of
\emph{odd clique chains} (block graphs whose block tree is a path and whose blocks are odd
complete graphs) we prove the precise refinements
\[
\pn(G)\le \floor{n/2}\ \ (b\ge 2), \qquad \pn(G)\le \floor{n/2}-1\ \ (b\ge 4),
\]
where $b$ is the number of blocks, along with the exact value $\pn(G)=\ceil{n/2}$ for a single
block. The local engine, of independent interest, is: \emph{for every odd $j$ and every
Hamilton path $P$ of $K_{2j+1}$, the graph $K_{2j+1}-P$ decomposes into $j$ paths.} The base
zigzag objects are classical \cite{LP92,Alspach08}; the statement for an \emph{arbitrarily
prescribed} Hamilton path appears to be new (the Oberwolfach problem and its variants concern
$2$-regular factors and do not cover path factors), and we could not find it in the Hamilton
decomposition literature.
\item \textbf{Fan-type theorems beyond bounded even-degree} (Section~\ref{sec:astar}). The Fan
line \cite{Fan05,BS21,FHZ23} has kept $\Delta(\EV(G))\le 3$ for twenty years. We prove three
class theorems that allow high vertices: Theorem~\ref{thm:Astar} (if $\EV(G)$ is triangle-free,
high E-degrees are compensated by an odd-degree condition, and an explicit \emph{alignment
condition} (A4) holds at one high vertex, then $\pn(G)\le\floor{n/2}$), Theorem~\ref{thm:sep}
(the same bound with no alignment hypothesis whenever distinct high vertices have no common
EV-neighbour), plus Remark~\ref{rem:CO}, which weakens separation to a parity condition on the
profiles of common EV-neighbours, and Theorem~\ref{thm:twohigh}, which settles the case of
exactly two high vertices that share EV-neighbours whenever some common EV-neighbour is clean
(E-degree two) or one high vertex has a private EV-neighbour of E-degree one;
Corollary~\ref{cor:twohigh} reduces the remaining two-high graphs to a corner case subject to
five simultaneous constraints at every common EV-neighbour of the two high vertices, of which
no computational instance was found; a band-free tightness theorem (Remark~\ref{rem:gap}) now
proves this corner empty on the constructed family proper whenever the counting gap is at most
two. The general statement is posed as
Conjecture~\ref{conj:Astar}; Proposition~\ref{prop:residue} pins down the residue class: a
minimal counterexample must have, at \emph{every} high vertex, a far high ``bad partner'' and,
except for the corner cases of Corollary~\ref{cor:twohigh}, at least three high vertices;
Remark~\ref{rem:gap} exhibits the concrete $23$-vertex double-star example on which the natural
induction stops at exactly this configuration. No counterexample is known: exact dynamic programs verify the conjecture with
zero violations on $579$ instances ($166$ from the computational study behind this paper, a further
$413$ in Appendix~\ref{app:comp}), and $15$ large verified decompositions achieve the bound.
Odd-order graphs in the (A1) class admit a density ceiling
(Proposition~\ref{prop:odddensity}), which in particular shows that the sharpness family of
Proposition~\ref{prop:oscsharp} lies exactly on the other side of (A1). Our witness family
$W(m)$, $m$ odd $\ge 5$, satisfies all hypotheses with $\EV(W(m))\cong K_{2,m}$, has treewidth
$5$, degeneracy $5$, and maximum degree at least $m+1$, hence escapes every published class we are aware
of --- in particular \cite{ZLH24,LY24,BP19,BCLS20,CFL} --- and
Proposition~\ref{prop:escape} shows that Theorem~\ref{thm:Astar} is currently the only known
result implying the $\floor{n/2}$ bound on $W(m)$. We also exhibit the odd semi-cliques
$K_{2k+1}-(k-1)K_2$, on which $\pn=\ceil{n/2}>\floor{n/2}$ for all $k$, showing that the
triangle-free requirement cannot simply be dropped.
\item \textbf{A conditional full chain: Gallai modulo explicit hypotheses}
(Section~\ref{sec:conditional}). We develop an endpoint-budget framework: the per-vertex
\emph{endpoint budget} $\beta(G,x)$ (Definition~\ref{def:endpointbudget}) and the budget
conjecture~$K$ that $\beta(G,x)\le\ceil{n/2}$ always. On the proved side we establish (i) an
absorption-cost lemma (Lemma~\ref{lem:T1}) that expresses exactly the cost of reinserting a set
of edges into a decomposition and unifies three obstruction phenomena from the literature;
(ii) that $K$ implies Gallai's conjecture (Theorem~\ref{thm:Kgallai}), so $K$ is a genuine
strengthening rather than a cheap consequence; and (iii) a reduction chain
(Theorem~\ref{thm:reduction}) assembling
$\text{Gallai}\Leftarrow K\Leftarrow\text{degree-two layer of }K\Leftarrow
\text{SC-PH}\Leftarrow\text{GAP-20a}^{\prime}\Leftarrow\text{(E)+(R)}$. The chapter's central
deliverable is an \emph{explicit hypothesis ledger}: each edge of the chain is stated in the
form ``Gallai modulo $\Pi$'' with its premise set $\Pi$ listed item by item --- including the
all-odd support case resting on the external reference \cite{Lo68} (flagged unverified against
the journal record for this use), the layer-~$N(k\ge4)$ status line (theorem at $k\le3$,
computational at $k=4,5$, open at $k\ge6$), and the open BP-odd conjecture --- so that every
conditional claim in the paper is checkable against exactly the premises it needs (five of the
eight premise conjuncts judgeable from the paper's own statements; the three bridge conjuncts
carried at registered pointer level, their definitions living in the campaign archive ---
Section~\ref{sec:conditional}).
\item \textbf{The machine-verification campaign package} (Sections~\ref{sec:machine},
\ref{sec:attack} and \ref{sec:lean}). The computational campaign behind the candidate
Proposition~\ref{prop:2663cand} (the so1 overdraw programme) was closed by a two-wave
adversarial adjudication (twelve review panels plus settlement, rework and errata executors;
zero red-level verdicts) with a registered three-tier outcome, restated as the master statement
of Section~\ref{sec:machine}: machine-theorem-level closure on the explicitly tested domains; a
paper-layer candidate; and zero claims beyond the tested domains. Section~\ref{sec:attack}
registers the current frontier of the companion selection-rule programme at $k=8$: the
candidate lemma G2-SL over the true quantifier domain $8{,}616=1{,}088+240+7{,}288$, which
passed a merged final adjudication with notes and stands as a \emph{candidate, not a theorem}.
Section~\ref{sec:lean} reports the Lean formalization base and the four-domain closure of the
hole-chain termination Theorem~D, whose four domain counts
$588/30{,}760/634{,}960/7{,}798{,}903$ --- each carried as a five-figure five-tuple --- are
directly machine-certified, together with the honest
residual-debt list. Six no-claim constraints bind this whole package (paragraph below);
the eighteen registered residual openings are collected in Section~\ref{sec:open}.
\end{enumerate}

\paragraph{Verification methodology.}
All exact values and all decompositions cited in this paper were produced and re-checked by an
independent machine pipeline (bitmask dynamic program plus a decomposition verifier asserting
that the parts are simple paths and exactly partition the edge set); scripts and certificate
files accompany the paper. Appendix~\ref{app:comp} lists, for each theorem, the verification
status and the machine anchors. Parts of the mathematical and bibliographic work were carried
out with computer assistance, including automated literature cross-checking against the
published record; the division of labour is summarized in Appendix~\ref{app:comp}.

\paragraph{Claim-strength declaration and the six no-claim constraints.}
The computational campaign behind Proposition~\ref{prop:2663cand} and its
2b dichotomy programme (Conjecture~\ref{conj:tdich}) was closed by a two-wave red/blue adjudication
(twelve review panels, with settlement, rework and errata executors; no
red-level verdict) with a three-tier outcome --- restated in full as the master statement of
Section~\ref{sec:machine}: machine-theorem-level closure
on the explicitly tested domains; a paper-layer candidate (fully proved
components plus one reworked patch whose independent final audit has
since landed and passed); and zero claims on the extrapolation domains. Six no-claim
constraints bind the whole account, including this paper: we do \emph{not} claim
(i) closure at $L\ge13$; (ii) general-$N$ paper closure; (iii) that the
four-piece (E)+(R) landing is complete; (iv) GL-C paper closure; (v) full
2b-side closure; (vi) theorem-level closure of
Proposition~\ref{prop:2663cand} (the campaign-internal labels are the
adjudication's; the verbatim register is carried in Appendix~\ref{app:errata}).
The eighteen residual openings are registered in
Appendix~\ref{app:finalverdict-open} and collected in Section~\ref{sec:open}.

\paragraph{Open problems.}
Section~\ref{sec:open} collects open problems: the even columns of Case~4 of \cite{CE18};
the removal of the structural conditions of Section~\ref{sec:astar}
(Conjecture~\ref{conj:Astar}); the even-$j$ residual of the Walecki engine; the
pair-hosting and density conjectures of Section~\ref{sec:conditional}; and the $4$-degenerate
residual class.

% ============================================================================
% §2 NOTATION AND PRELIMINARIES. Migration anchor: transplanted VERBATIM from
% main.tex :262-433 (source §2 Preliminaries, incl. Lemmas LB1/LB2/sum,
% Theorems lovasz/pyber/fan, Fan reinsertion lemmas with transcription and
% lane_154 downgrade notes, odd semi-clique definition, Walecki partition) —
% see ANCHOR comment at the end of the transplanted block. New subsection
% \ref{subsec:machnotation} (unified campaign notation) added per P-ARCH §2.2
% item 3; T_res definition written with f in {6,7} per erratum N-COLD-1
% (FRONTIER.md REVISE-29), t := |V(T)| = f+1 per N-COLD-2.
% ============================================================================
\section{Preliminaries}\label{sec:prelim}

All graphs are finite, simple and undirected. A \emph{path} has at least one edge and no
repeated vertex; a \emph{path decomposition} of $G$ is a set $D$ of paths whose edge sets
partition $E(G)$; $\pn(G)$ is the minimum size of a path decomposition. For a decomposition $D$
and a vertex $x$ we write $D(x)$ for the number of paths of $D$ having $x$ as an endpoint. A
\emph{block} of $G$ is a maximal $2$-connected subgraph or a bridge (together with its
endpoints); every two blocks share at most one vertex, a \emph{cut vertex}. $G$ is
\emph{$k$-degenerate} if every subgraph has a vertex of degree at most $k$. The
\emph{even-degree subgraph} $\EV(G)$ is the subgraph induced by the vertices of even degree;
for $v\in V(G)$ we write $\dEV(v)$ for the degree of $v$ inside $\EV(G)$, the \emph{E-degree}
of $v$. Note that $\dEV(v)\ge 4$ implies in particular that $v$ has even degree.

We collect the standard tools. Proofs of the first three are included for completeness; they
are folklore.

\begin{lemma}[odd-degree lower bound; endpoint parity]\label{lem:LB1}
Let $D$ be a path decomposition of $G$. Then:
\begin{enumerate}[label=(\roman*),leftmargin=2em]
\item $2|D|\ge \mathrm{odd}(G)$, where $\mathrm{odd}(G)$ is the number of odd-degree vertices;
in particular $\pn(G)\ge \mathrm{odd}(G)/2$;
\item every odd-degree vertex is an endpoint of at least one path of $D$;
\item for every vertex $v$, $D(v)\equiv \deg_G(v)\pmod 2$.
\end{enumerate}
\end{lemma}

\begin{proof}
For a path $P$ and a vertex $v$, the degree of $v$ in $P$ is $0$, $1$ or $2$, and it is $1$
exactly when $v$ is an endpoint of $P$. Summing over $D$,
$\deg_G(v)=\sum_{P\in D}\deg_P(v)\equiv D(v)\pmod 2$, which gives (iii); (ii) is (iii) for odd
$\deg_G(v)$; and (i) follows by charging each odd-degree vertex to one path ending at it, each
path being charged at most twice.
\end{proof}

\begin{lemma}[edge-count lower bound]\label{lem:LB2}
If $G$ is connected with $n\ge 2$ vertices and $m$ edges, then
$\pn(G)\ge \ceil{m/(n-1)}$.
\end{lemma}

\begin{proof}
Each path has at most $n-1$ edges, so $k$ paths cover at most $k(n-1)$ edges.
\end{proof}

\begin{lemma}[disconnected sum]\label{lem:sum}
If $G$ has components $G_1,\dots,G_c$ on $n_1,\dots,n_c$ vertices, then
$\sum_i \floor{n_i/2}\le \floor{n/2}$.
\end{lemma}

\begin{proof}
Immediate from $\sum_i \floor{n_i/2}\le \floor{\sum_i n_i /2}$ for integers.
\end{proof}

\begin{theorem}[Lov\'asz \cite{Lo68}]\label{thm:lovasz}
Every $n$-vertex graph (not necessarily connected) has a decomposition of its edge set into at
most $\floor{n/2}$ paths and cycles.
\end{theorem}

\begin{theorem}[Pyber \cite{Pyb96}]\label{thm:pyber}
If $\EV(G)$ is a forest, then $\pn(G)\le \floor{n/2}$.
\end{theorem}

\begin{theorem}[Fan \cite{Fan05}]\label{thm:fan}
If every block of $\EV(G)$ is triangle-free with maximum degree at most $3$, then
$\pn(G)\le \floor{n/2}$.
\end{theorem}

Our proofs in Section~\ref{sec:astar} use Fan's reinsertion engine, whose two lemmas we state
in the form in which they are restated in \cite{BS21} (Lemmas 6 and 7 there; in \cite{Fan05}
they appear as Lemmas 3.4 and 3.6). For a decomposition $D$ and a vertex $x$, call a neighbour
$w$ of $x$ \emph{active} if $D(w)>0$.

\begin{lemma}[addible edge; {\cite{Fan05},\cite{BS21}}]\label{lem:fan6}
Let $uv\in E(G)$ and let $D'$ be a path decomposition of $G-uv$. If
$D'(v)>\ |\{w\in N_{G-uv}(u) : D'(w)=0\}|$, then $G$ has a path decomposition $D$ with
$|D|=|D'|$, $D(u)=D'(u)+1$, $D(v)=D'(v)-1$, and $D(w)=D'(w)$ for all other $w$. One says that
$uv$ is \emph{addible towards $u$}.
\end{lemma}

\emph{Transcription and verification note.} The proof of Lemma~\ref{lem:fan6} in \cite{Fan05}
could not be consulted; the statement follows the restatement in \cite{BS21} (Lemma 6) and has
not been independently checked against the original. A strengthening of the lemma that
additionally requires the underlying re-arrangement to preserve the endpoint profile of $D'$
pointwise is false in general (a five-vertex counterexample is machine-certified); the lemma as
stated imposes no such requirement and has no known counterexample. Its machine evidence is
direction-sensitive: the historical $n=6$ scan of $850{,}570{,}575$ hypothesis-satisfying
instances tested one fixed orientation per edge, whereas the applications below use both
orientations; the direction-corrected bidirectional rescan is complete on $n\le6$: all
$482{,}080$ ordered instances at $n\le5$ ($241{,}040$ per orientation) and all
$1{,}701{,}141{,}150$ ordered instances at $n=6$ ($850{,}570{,}575$ per orientation, the
fixed-orientation half exactly reproducing the historical census), with zero violations
(Appendix~\ref{app:comp}).

\emph{Downgrade note (lane\_154).} The Lean formalization of this lemma uses the four-clause
form obtained by deleting the final pointwise clause $D(w)=D'(w)$ for the remaining vertices:
no application below (including the one inside Lemma~\ref{lem:fan7}) consumes that clause, the
formal re-assembly route (a rearrangement preserving the two-point profile at $u,v$ plus the
extension step) proves exactly the four-clause form, and the pointwise clause is retained only
as a pending existence statement. The semantics of every application below are unaffected.

\begin{lemma}[star absorption; {\cite{Fan05},\cite{BS21}}]\label{lem:fan7}
Let $B=\{uw_1,\dots,uw_h\}$ be a set of $h$ edges incident to $u$, let $D'$ be a path
decomposition of $G-B$ in which every neighbour of $u$ \emph{in $G$} has $D'\ge 1$, and let
$uw$ be a specified edge of $B$. Then there is a set $A\subseteq B$ with $uw\in A$ and
$|A|\ge\ceil{h/2}$ such that the edges of $A$ can be added back one by one, each by
Lemma~\ref{lem:fan6} towards $u$, preserving $|D|$ and increasing $D(u)$ by one at each step.
\end{lemma}

\emph{Reading note.} The hypothesis of Lemma~\ref{lem:fan7} is read over the neighbours of $u$
\emph{in $G$}, so that it includes the star vertices $w_i$ themselves, which are not
neighbours of $u$ in $G-B$; this stronger reading is the one verified in every application
below (the proof of Theorem~\ref{thm:Astar}, Step~3, makes it explicit). As with
Lemma~\ref{lem:fan6}, the proof in \cite{Fan05} (Lemma 3.6) could not be consulted, and the
transcription from the restatement in \cite{BS21} (Lemma 7) has not been independently checked
against the original.

\begin{definition}[odd semi-clique; {\cite{BP19}}]\label{def:osc}
An \emph{odd semi-clique} is a graph obtained from the complete graph on $2k+1$ vertices by
deleting at most $k-1$ edges.
\end{definition}

\begin{observation}[density obstruction]\label{obs:density}
An odd semi-clique on $n=2k+1$ vertices has
$m\ge \binom{2k+1}{2}-(k-1)=2k^2+1$ edges, so by Lemma~\ref{lem:LB2}
$\pn\ge \ceil{(2k^2+1)/2k}=k+1=\ceil{n/2}$; in particular no odd semi-clique admits a
$\floor{n/2}$-path decomposition.
\end{observation}

Finally we record the classical Walecki partition of the complete graph of odd order, in the
form used in Section~\ref{sec:chains}. Fix $j\ge 1$ and let
$V=\{\infty\}\cup \mathbb{Z}_{2j}$. For $k\in\mathbb{Z}_{2j}$ let $C_k$ be the cycle
\[
C_k \;=\; \infty,\, k,\, k+1,\, k-1,\, k+2,\, k-2,\;\dots,\, k+(j-1),\, k-(j-1),\, k+j,\, \infty
\]
(indices mod $2j$); equivalently the finite vertices of $C_k$ are $p_0=k$ and
$p_s=k+(-1)^{s+1}\ceil{s/2}$ for $s=1,\dots,2j-1$.

\begin{lemma}[Walecki partition; classical, cf.\
\cite{LP92,Alspach08}]\label{lem:walecki}
For $k=0,1,\dots,j-1$ the cycles $C_k$ are edge-disjoint and
$\bigsqcup_{k=0}^{j-1}E(C_k)=E(K_{2j+1})$.
\end{lemma}

\begin{proof}
Each $C_k$ is a cycle: the finite vertices $k, k\pm1,\dots,k\pm(j-1),k+j$ are $2j$ distinct
elements covering $\mathbb{Z}_{2j}$. We show every edge lies in exactly one $C_k$, $k$ taken
mod $j$. Write $\delta(x,y)$ for the circular distance of $x,y\in\mathbb{Z}_{2j}$, i.e.\ the
minimum of $\pm(x-y)$ mod $2j$ in $\{0,\dots,j\}$.

Reading off consecutive differences in the sequence $p_s$ (as above, $p_{s+1}-p_s=\pm(s+1)$),
one checks that the finite edges of $C_k$ are exactly the pairs
$\{k+t,\,k-t\}$ for $t\in\mathbb{Z}_{2j}\setminus\{0,j\}$ and the pairs
$\{k+t,\,k-t+1\}$ for $t\in\mathbb{Z}_{2j}$, together with the two $\infty$-edges
$\{\infty,k\}$, $\{\infty,k+j\}$. (In each of the two forms an edge is produced twice, once
from each endpoint, so it is counted once; $t=0$ and $t=j$ would give loops in the first form
and are excluded.) Hence $C_k$ has, among finite edges, two edges of each circular distance
$1,\dots,j-1$ and one edge of distance $j$.

Now fix an edge $e=\{x,y\}$ of $K_{2j+1}$ and let $d=\delta(x,y)$.
\begin{itemize}[leftmargin=1.5em]
\item If $y=x+j$ (the diameter case, $d=j$): for $j$ odd, $e$ has the form $\{k+t,k-t+1\}$,
which forces $2k\equiv 2x+d-1\pmod{2j}$, a unique solution mod $j$ (the right side is even);
for $j$ even, $e$ has the form $\{k+t,k-t\}$ with $2t\equiv d=j$, forcing
$2k\equiv 2x+j\pmod{2j}$, again a unique solution mod $j$. So $e$ lies in exactly one $C_k$.
\item If $d<j$ is even: $e=\{x,x+d\}=\{k+t,k-t\}$ forces $2k\equiv 2x+d\pmod{2j}$, unique
mod $j$.
\item If $d<j$ is odd: $e=\{x,x+d\}=\{k+t,k-t+1\}$ forces $2k\equiv 2x+d-1\pmod{2j}$, unique
mod $j$.
\item If $e=\{\infty,x\}$: $e\in C_k$ iff $x\in\{k,k+j\}$, i.e.\ $k\equiv x$ or $x+j$
(mod $2j$), a unique class mod $j$.
\end{itemize}
In each case $e$ lies in exactly one of $C_0,\dots,C_{j-1}$, and each $C_k$ has $2j+1$ edges,
so the $j(2j+1)=\binom{2j+1}{2}$ edges are partitioned.
\end{proof}
% !ANCHOR main.tex:L262-L433 -> main_submission.tex §2 (verbatim transplant,
% body complete; label sec:prelim and all inner labels preserved).

\subsection{Notation for the conditional and machine chapters}\label{subsec:machnotation}

We collect, in one place, the notation used by the campaign chapters
(Sections~\ref{sec:machine}, \ref{sec:attack}, \ref{sec:conditional} and \ref{sec:lean}).

\begin{definition}[endpoint budget]\label{def:endpointbudget}
For a connected graph $G$ and an even-degree non-isolated vertex $x$ of $G$, the \emph{endpoint
budget} $\beta(G,x)$ is the minimum cardinality of a path decomposition $D$ of $G$ in which $x$
is an endpoint of at least two paths. The \emph{budget conjecture~$K$} is the statement that
$\beta(G,x)\le\ceil{n/2}$ for every connected $G$ and every even-degree non-isolated vertex $x$;
the formal definition and the conjecture are restated in
Section~\ref{sec:conditional} (Definition~\ref{def:beta}, Conjecture~\ref{conj:K}).
\end{definition}

\begin{sloppypar}
\emph{Difference-board instances at $k=3$}
(campaign ledgers \texttt{lane\_260\_muk3.md} and \texttt{lane\_266\_muk3final.md}; domain
language transcribed from the source account, see Appendix~\ref{app:mu2663}). A
\emph{$k=3$ instance} is a triple $(P,T,d)$ with $P\subseteq\mathbb{Z}_N$, $p=|P|$,
three pairwise
distinct anchor offsets $T=(t_1,t_2,t_3)$, and demands $d=(d_1,d_2,d_3)$ meeting
(H1) $d_i\le p-[t_i\in P]$ for each $i$ and (H2) $\beta\ge 1$, where $n=d_1+d_2+d_3$ and
$\beta:=p-n$. Writing $\delta_{ij}:=t_i-t_j$, a \emph{legal step pair} is a choice of
$C_1\subseteq P\setminus\{t_1\}$, $|C_1|=d_1$, and
\[
  C_2\subseteq P_2(C_1):=P\setminus\bigl(C_1\cup((C_1-\delta_{12})\cap P)\cup
  (\{t_2\}\cap P)\bigr),\qquad |C_2|=d_2;
\]
the \emph{three-source account} of the step is
\[
\begin{aligned}
  g_{13}^{+}&=\bigl|(C_1-\delta_{13})\cap P\setminus(C_1\cup C_2)\bigr|,\\
  g_{23}^{+}&=\bigl|(C_2-\delta_{23})\cap P\setminus(C_1\cup C_2\cup
    ((C_1-\delta_{13})\cap P))\bigr|,\\
  e_3&=\bigl|\{t_3\}\cap P\setminus(C_1\cup C_2)\bigr|,
\end{aligned}
\]
with $g_{13}^{+}+g_{23}^{+}+e_3=:\mathrm{val}$; the instance's
$\min\mathrm{val}$ is the minimum of $\mathrm{val}$ over all legal step pairs.
The \emph{so1 (aligned) family} is the subfamily
$\delta_{13}\equiv k\,\delta_{12}\pmod N$ with $g=\gcd(\delta_{12},N)$,
$L=\mathrm{ord}(\delta_{12})$, $k\in[2,L)$ (the anchors being pairwise distinct); on each
residue chain of $\langle\delta_{12}\rangle$ one uses main-chain coordinates
$t_1=0$, $t_2=-1$, $t_3=-k$ in $\mathbb{Z}_L$, and a \emph{$k$-cycle} is an orbit of
$\langle-k\,\delta_{12}\rangle$ (length $L_c\mid L$).
\end{sloppypar}

\begin{definition}[residual tree types at $k=8$; corrected quantifier
domain]\label{def:tres}
At $k=8$ one works over $R=\mathbb{Z}_{16}\sqcup\{\mathrm{Inf}\}$ with a class function
$r$ on the pairs of $R$. We write $\mathcal{T}_{\mathrm{res}}$ for the family of the
$30$ \emph{residual tree types} with $f\in\{6,7\}$ edge classes --- $9$ types at $f=6$ and
$21$ types at $f=7$ --- independently regenerated by both audit reports. For
$T\in\mathcal{T}_{\mathrm{res}}$ we write
$t:=|V(T)|=f+1$ for the vertex count of $T$, and $\Delta(T)$ for the class-count parameter
of $T$. For a free set $S\subseteq R$ with $|S|=n$ and a centre $c\in S$,
$\mathrm{cap}_S(c)=|\{r(\{c,y\}):y\in S\setminus\{c\}\}|$ (taken over centres
\emph{actually present in} $S$, the F-C4-1 convention) and
$\mathrm{maxcap\_avail}(S)=\max_{c\in S}\mathrm{cap}_S(c)$. The \emph{slack} of $S$
(relative to $T$) is $\mathrm{slack}:=n-t$.
\end{definition}

\begin{remark}[quantifier domains; erratum N-COLD-1/N-COLD-2 carried on the
notation]\label{rem:tresdomain}
The selection-rule candidate G2-SL of Section~\ref{sec:attack} is quantified over
$T\in\mathcal{T}_{\mathrm{res}}$ and free $S\subseteq R$ with $\mathrm{slack}\ge1$ (i.e.\
$n\ge t+1$) and $\mathrm{maxcap\_avail}(S)=\Delta(T)$ (guards G2-a and G2-b), with
$\Delta\in\{3,4,5,6\}$ at $k=8$; the slack-$0$ layer is false (three T19 sets at $n=8$).
Two erratum notes are carried on the notation itself, per the campaign frontier register:
(i) the type family is \emph{defined} with $f\in\{6,7\}$ --- a literal reading ``$f\le7$''
would admit $f\le5$ types not present in any machine account (a $35$-type literal domain of
$8{,}672$ instances against the audited $8{,}616$), so the definition above fixes
$f\in\{6,7\}$; (ii) the vertex count $t:=|V(T)|=f+1$ is part of the notation, not left to
context.
\end{remark}
% !ANCHOR notation sources: main.tex :3918-3938 (k=3 instance, H1/H2, val),
% :4263-4273 (so1 family, main-chain coordinates, k-cycles),
% :6629-6658 (R, 30 residual types, t=f+1, slack, cap_S, maxcap_avail,
% guards G2-a/G2-b, Delta range) — transcribed with N-COLD-1 f∈{6,7} fix.

% ============================================================================
% §3 FAN-TYPE THEOREMS — VERBATIM TRANSPLANT (construction batch 4, deep_P17).
% Source: frozen ledger main.tex :434-1020 (P-ARCH migration row "§3 (434-1020)
% 保留 -> 新 §3 原文不动"; 587 source lines relocated byte-for-byte). This closes
% the eleven §3-owned dangling references registered by deep_P14/P16 (thm:Astar,
% thm:sep, thm:twohigh, cor:twohigh, rem:CO, rem:gap, conj:Astar, prop:residue,
% prop:escape, prop:odddensity, prop:oscsharp): all eleven labels are defined in
% this block (source lines :451-:975).
% ============================================================================
\section{Fan-type theorems beyond bounded even-degree}\label{sec:astar}

Throughout this section $G$ is a connected graph, $n=|V(G)|$, and $\dEV(v)$ denotes the E-degree
of $v$. Recall that a vertex with $\dEV\ge 4$ has even degree in $G$; call such a vertex
\emph{high}, and call a vertex with $1\le \dEV\le 3$ \emph{low}. All published results in the
Fan line \cite{Fan05,BS21,FHZ23} require the blocks of $\EV(G)$ to have maximum degree at most
$3$, i.e.\ they require that \emph{no} vertex be high. Theorem~\ref{thm:Astar} allows high
vertices, under an odd-degree compensation condition (A2) together with an alignment condition
(A4) at one high vertex; Theorem~\ref{thm:sep} removes the alignment condition when distinct
high vertices have no common EV-neighbour, and Remark~\ref{rem:CO} weakens separation to a
parity condition on the profiles of common EV-neighbours; Theorem~\ref{thm:twohigh} treats the
case of exactly two high vertices that share EV-neighbours, and Corollary~\ref{cor:twohigh}
collects the at-most-two-high layer, leaving only a corner open there, subject to five
simultaneous constraints at every common EV-neighbour of the two high vertices; the
band-free tightness theorem (Remark~\ref{rem:gap}, Appendix~\ref{app:band}) empties this
corner on the constructed family proper whenever the counting gap is at most two.

\begin{theorem}[alignment at a high vertex]\label{thm:Astar}
Let $G$ be a connected graph such that
\begin{enumerate}[label={},leftmargin=2.6em]
\item[(A1)] $\EV(G)$ is triangle-free;
\item[(A2)] every high vertex $v$ has odd $\dEV(v)$, and every EV-neighbour $w$ of $v$
satisfies $\dEV(w)\le 3$;
\item[(A4)] there is a high vertex $u$ such that every even vertex
$x\notin\{u\}\cup N_{\EV}(u)$ satisfies
$\dEV(x)-|N_{\EV}(x)\cap N_{\EV}(u)|\le 3$, where
$N_{\EV}(x)$ denotes the set of EV-neighbours of $x$.
\end{enumerate}
Then $\pn(G)\le \floor{n/2}$.
\end{theorem}

Here $N_{\EV}(x)$ denotes the set of EV-neighbours of $x$, i.e.\ its neighbours inside
$\EV(G)$.\footnote{The numbering skips (A3): a condition so numbered in an earlier draft of
this section was absorbed into the separation hypothesis of
Theorem~\ref{thm:sep}, and the remaining labels were kept unchanged.}

\begin{proof}
Let $F=\{uv : v\in N_{\EV}(u)\}$ be the $E$-star at $u$, $|F|=k$, and let $\tilde G=G-F$. We
first analyse the parity and structure of $\tilde G$.

\emph{Step 1 (structure of $\tilde G$).}
Every edge of $F$ joins two even-degree vertices of $G$, so odd-degree vertices of $G$ lose no
edges and stay odd. The vertex $u$ loses $k$ edges, an odd number, so $u$ is odd in
$\tilde G$. Each $v\in N_{\EV}(u)$ loses exactly one edge, so $v$ is odd in $\tilde G$.
Every other vertex keeps its degree. Hence
$V(\EV(\tilde G))=\{x\notin\{u\}\cup N_{\EV}(u) : \deg_G(x)\text{ even}\}$, and since
$\EV(\tilde G)$ is an induced subgraph of $\EV(G)$, it is triangle-free by (A1).

Moreover every $x\in V(\EV(\tilde G))$ has $\dEV$ in $\tilde G$ equal to
$\dEV(x)-|N_{\EV}(x)\cap N_{\EV}(u)|\le 3$ by (A4). Thus \emph{every} vertex of $\EV(\tilde G)$
has E-degree at most $3$, so every block of $\EV(\tilde G)$ is triangle-free of maximum degree
at most $3$. Applying Theorem~\ref{thm:fan} to each component of $\tilde G$ (the hypotheses are
component-wise) and summing with Lemma~\ref{lem:sum},
\begin{equation}\label{eq:tildebound}
\pn(\tilde G)\;\le\;\textstyle\sum_c \floor{n_c/2}\;\le\;\floor{n/2},
\end{equation}
where $n_c$ are the component orders. Fix a decomposition $\tilde D$ of $\tilde G$ with
$|\tilde D|\le\floor{n/2}$.

\emph{Step 2 (hosts).}
The neighbours of $u$ in $\tilde G$ are precisely its odd-degree neighbours in $G$; they are
untouched and odd in $\tilde G$, so by Lemma~\ref{lem:LB1}(ii) each has $\tilde D\ge 1$.

\emph{Step 3 (reinsertion of the addible part).}
Apply Lemma~\ref{lem:fan7} with $B=F$, $h=k$, and specified edge $uv_1$ (any fixed element of
$F$). The hypothesis of Lemma~\ref{lem:fan7} is satisfied: every neighbour of $u$ has
$\tilde D\ge 1$ by Step 2 (this holds whether the hypothesis is read over the neighbours of $u$
in $\tilde G$ or in $G$, since the star vertices $v\in N_{\EV}(u)$ are odd in $\tilde G$, each
having lost a single edge, and have $\tilde D\ge 1$ by Lemma~\ref{lem:LB1}(ii)). We obtain
$A\subseteq F$ with $|A|\ge\ceil{k/2}$ which can be added back towards $u$, producing a
decomposition $D_1$ of $\tilde G\cup A$ with $|D_1|=|\tilde D|$ and
$D_1(u)\ge \tilde D(u)+\ceil{k/2}\ge 1+\ceil{k/2}$, the bound $\tilde D(u)\ge1$ coming from
Lemma~\ref{lem:LB1}(ii) applied to $u$, which is odd in $\tilde G$.

\emph{Step 4 (reinsertion of the remainder).}
Let $L=F\setminus A$; $|L|=\floor{k/2}=(k-1)/2$. Process the edges of $L$ one at a time.
Suppose $j-1$ edges of $L$ have been reinserted, giving a decomposition $D_j$ of
$\tilde G\cup A\cup(\text{first } j-1 \text{ edges of } L)$ with
\[
D_j(u)\;\ge\;1+\ceil{k/2}-(j-1).
\]
Let $uv\in L$ be the next edge. We apply Lemma~\ref{lem:fan6} \emph{towards $v$}, i.e.\ with
the roles of $u$ and $v$ in the statement exchanged: the condition required is
\[
D_j(u)\;>\;\bigl|\{w\in N_{G_j}(v) : D_j(w)=0\}\bigr|,
\]
where $G_j$ is the current graph, and the operation then yields a decomposition $D_{j+1}$ of
$G_j+uv$ with $D_{j+1}(v)=D_j(v)+1$, $D_{j+1}(u)=D_j(u)-1$, and $|D_{j+1}|=|D_j|$.

We bound the right-hand side. If $D_j(w)=0$ then $w$ has even degree in $G_j$ by
Lemma~\ref{lem:LB1}(iii). We classify the neighbours of $v$ and their endpoint counts.
Odd-degree vertices of $G$ are never incident with a star edge, keep their degree, and
remain odd throughout; they always have $D\ge1$. A star vertex $v'\in N_{\EV}(u)$,
$v'\ne v$, is \emph{not a neighbour of $v$ at all}: the vertices $v$ and $v'$ are both
even-degree vertices of $G$, so the edge $vv'$ would close the triangle $u,v,v'$ in
$\EV(G)$, contradicting (A1). (This is just as well: a star vertex whose edge has been
reabsorbed towards $u$ by Lemma~\ref{lem:fan7} may well have $D=0$, each absorption step
lowering $D(v')$ by one.) The vertex $u$ \emph{is} a neighbour of $v$ and has
regained at least $\ceil{k/2}\ge3$ edges, so $D_j(u)\ge3>0$. The only remaining
neighbours of $v$ that were even in $G$ and lie outside $\{u\}\cup N_{\EV}(u)$ are the
members of $N_{\EV}(v)\setminus\{u\}$, of which there are $\dEV(v)-1\le2$. Hence
\[
\bigl|\{w\in N_{G_j}(v) : D_j(w)=0\}\bigr|
\;\le\; \bigl|\{w\in N_{\EV}(v) : w\ne u\}\bigr|
\;=\;\dEV(v)-1\;\le\;2,
\]
using that $v$ is low, i.e.\ $\dEV(v)\le 3$ by (A2). At the last
step $j=(k-1)/2$ we have
$D_j(u)\ge 1+\ceil{k/2}-(j-1)=1+(k+1)/2-(k-3)/2=3>2$, and for earlier steps the budget is only
larger. So each edge of $L$ is addible towards its star vertex $v$, free of charge.

After all $k$ edges are reinserted we obtain a decomposition of $G$ with
$|D|=|\tilde D|\le\floor{n/2}$, completing the proof.
\end{proof}

The alignment condition is not the only route through a high vertex; the following theorem
needs no alignment hypothesis at all, at the price of a separation condition between high
vertices.

\begin{theorem}[separated high vertices]\label{thm:sep}
Assume {\rm(A1)}, {\rm(A2)}, and moreover that any two distinct high vertices of
$G$ have no common EV-neighbour. Then $\pn(G)\le\floor{n/2}$.
\end{theorem}
\begin{proof}
We prove the stronger statement, for all graphs in the class connected or
not, by strong induction on $|E(G)|$. If there is no high vertex, every E-degree is at
most $3$ and Theorem~\ref{thm:fan} applies component-wise, so
Lemma~\ref{lem:sum} gives the bound. Otherwise pick a high vertex $u$ and let
$k=\dEV(u)\ge 4$, which is odd by {\rm(A2)}; let $F$ be the E-star at $u$ and
$\tilde G=G-F$. Every star vertex has odd degree in $\tilde G$ (it loses the
single edge $uv$), and $u$ has odd degree in $\tilde G$; the even-degree
vertices of $\tilde G$ are precisely the even-degree vertices of $G$ outside
$\{u\}\cup N_{\EV}(u)$, and for such a vertex $x$ its E-degree in $\tilde G$
equals $\dEV(x)-|N_{\EV}(x)\cap N_{\EV}(u)|$; when $x$ is high the subtraction
term is $0$, because distinct high vertices have no common EV-neighbour. Hence $\tilde G$
again satisfies {\rm(A1)}, {\rm(A2)} and the separation condition, so the
induction hypothesis gives a decomposition $\tilde D$ of $\tilde G$ with
$|\tilde D|\le\floor{n/2}$. The neighbours of $u$ in $\tilde G$ are its
odd-degree neighbours in $G$, so each has $\tilde D\ge 1$
(Lemma~\ref{lem:LB1}(ii)); the star vertices are odd in $\tilde G$, each having
lost a single edge, so they too have $\tilde D\ge1$ by
Lemma~\ref{lem:LB1}(ii), and the hypothesis of Lemma~\ref{lem:fan7} holds read
over $G$; Lemma~\ref{lem:fan7} returns an addible set
$A\subseteq F$ with $|A|\ge\ceil{k/2}$ towards $u$, giving
$D_1(u)\ge 1+\ceil{k/2}$. Each edge of $L=F\setminus A$, $|L|=\floor{k/2}$, is
then reinserted towards its star vertex $v$ via Lemma~\ref{lem:fan6}: the
vertices $w$ with $D_j(w)=0$ are even in $G_j$, and as in the proof of
Theorem~\ref{thm:Astar} they are confined to $N_{\EV}(v)\setminus\{u\}$, of
size at most $\dEV(v)-1\le 2$; at the last step $D_j(u)\ge 3>2$. This rebuilds
a decomposition of $G$ with $\floor{n/2}$ paths.
\end{proof}

\begin{remark}\label{rem:CO}
The proof of Theorem~\ref{thm:sep} goes through, verbatim up to one parity
check, with separation replaced by the following weaker condition: for any two
distinct high vertices $x,z$ and any set $S$ of high vertices disjoint from
$\{x,z\}$, the number of common EV-neighbours $w$ of $x$ and $z$ with
$N_{\EV}(w)\cap H\cap S=\emptyset$ is even (equivalently: for each pair $x,z$,
the common EV-neighbours are split by their ``profile'' $\pi(w)=N_{\EV}(w)\cap H$
into classes of even size). Indeed this condition is again inherited after a
star deletion, because a vertex of $\EV(\tilde G)$ is precisely an even vertex
seeing no member of $\{u\}\cup N_{\EV}(u)$, so the counts for $\tilde G$ with
$S'$ are the original counts for $S=\{u\}\cup S'$.
\end{remark}

Separation and its profile relaxation still leave open the case of high vertices that do share
EV-neighbours. The next theorem settles the case of exactly two high vertices except for a
corner recorded afterwards.

\begin{theorem}[two high vertices]\label{thm:twohigh}
Assume {\rm(A1)} and {\rm(A2)}, let $u,x$ be the \emph{only} high vertices of $G$, and suppose
$t_{ux}:=|N_{\EV}(u)\cap N_{\EV}(x)|\ge1$. If either
\begin{enumerate}[label={\rm(\roman*)},leftmargin=2em]
\item some common EV-neighbour $v_1$ of $u$ and $x$ has $\dEV(v_1)=2$, or
\item $x$ has a private EV-neighbour $y^*$, i.e.\ $y^*\in N_{\EV}(x)\setminus N_{\EV}(u)$,
with $\dEV(y^*)=1$,
\end{enumerate}
then $\pn(G)\le\floor{n/2}$.
\end{theorem}

\begin{proof}
Let $F$ be the E-star at $u$, $k=\dEV(u)$, and delete one further edge $e^*$ at $x$:
$e^*=xv_1$ in~(i), $e^*=xy^*$ in~(ii), giving $\tilde G=G-F-e^*$.
In $\tilde G$ the vertices $u$ and $x$ are odd, and so is every former star vertex (each
lost the single edge $uv$); in~(ii) the far endpoint $y^*$ of $e^*$ is odd as well (it lost
$xy^*$). The remaining even vertices --- in~(i) the vertex $v_1$ (which lost both $uv_1$
and $xv_1$ and stays even) together with the far even vertices and the private
EV-neighbours of $x$; in~(ii) the far even vertices other than $y^*$ and the private
EV-neighbours of $x$ other than $y^*$ --- all have E-degree at most $3$ in $\tilde G$: the
private ones by
{\rm(A2)} at $x$, the far ones because $u$ and $x$ are the only high vertices, and $v_1$
because $\dEV(v_1)=2$. Hence $\EV(\tilde G)$ is triangle-free, being an induced subgraph of
$\EV(G)$, of maximum E-degree at most $3$, so $\tilde G$ lies in the class of
Theorem~\ref{thm:fan} and has, component-wise with Lemma~\ref{lem:sum}, a decomposition
$\tilde D$ with $|\tilde D|\le\floor{n/2}$. In~(i) reinsert $e^*$ towards $v_1$
(Lemma~\ref{lem:fan6}: $\tilde D(x)\ge1$ since $x$ is odd, while every neighbour of $v_1$ in
$\tilde G$ is odd except $x$ --- the even-degree neighbours of $v_1$ in $G$ are exactly $u$ and
$x$ --- so the passing count is $0$); this activates $v_1$. In~(ii) no activation is needed:
every star vertex has lost exactly one edge and is odd. Then Lemma~\ref{lem:fan7} (hypothesis
read over $G$, as in the proof of Theorem~\ref{thm:Astar}) returns $A\subseteq F$ with
$|A|\ge\ceil{k/2}$ addible towards $u$, and each edge of $L=F\setminus A$ is reinserted towards
its star vertex as in the proof of Theorem~\ref{thm:Astar} (passing count $\le\dEV(v)-1\le2$,
budget $\ge3$ at the last step). In~(ii) the edge $xy^*$ is reinserted towards $y^*$, whose
neighbours except $x$ are all odd, so this step is legal at any position of the sequence. All
$k+1$ deleted edges are restored with $|D|=|\tilde D|\le\floor{n/2}$.
\end{proof}

\begin{corollary}\label{cor:twohigh}
If $G$ satisfies {\rm(A1)}--{\rm(A2)} and has at most two high vertices, then
$\pn(G)\le\floor{n/2}$, unless $G$ has exactly two high vertices $u,x$ with $t_{ux}\ge1$,
every common EV-neighbour of $u$ and $x$ of E-degree $3$, no private EV-neighbour of E-degree
$1$ on either side, and such that, for \emph{every} common EV-neighbour $v$ of $u$ and $x$,
writing $z$ for the third EV-neighbour of $v$, $F$ for the E-star at $u$ and
$\tilde G=G-F-\{xv\}$, all of the following hold:
$\pn(\tilde G)=\floor{n/2}$; $\beta(\tilde G,x)=\floor{n/2}+1$ (cf.\
Section~\ref{sec:conditional}% !MIGRATION NOTE (deep_P17): source main.tex reads
% Section~
ef{sec:framework}; the conditional chapter is 
ef{sec:conditional}
% in the submission (P14 rename; same repair as the sec6-10 self-references).
); $z$ is buried (i.e.\ $D(z)=0$) in every optimal decomposition of
$\tilde G$; in every optimal decomposition $\tilde D$ of $\tilde G$, either $\tilde D(v)=0$,
or $\tilde D(v)\ge2$ and at least $\tilde D(v)$ private EV-neighbours of $x$ have
$\tilde D=0$; and $\tilde G$ has at most one component of odd order.
\end{corollary}
\begin{proof}
$t_{ux}=0$ is Theorem~\ref{thm:sep} (with only two high vertices the separation condition
reduces to $t_{ux}=0$). The two alternatives are Theorem~\ref{thm:twohigh}. For the remaining
graphs fix a common EV-neighbour $v$, write $z$ for its third EV-neighbour and
$\tilde G=G-F-\{xv\}$; then $\EV(\tilde G)$ again has maximum E-degree $\le3$, and each of the
following five routes closes the case for this $v$, the star at $u$ being absorbed afterwards
exactly as in the proof of Theorem~\ref{thm:twohigh}:
\begin{enumerate}[label={\rm(\alph*)},leftmargin=2em]
\item $\pn(\tilde G)\le\floor{n/2}-1$; take an optimal decomposition $\tilde D$. If
$\tilde D(x)\ge3$ we are on route~(b). Otherwise $\tilde D(x)=1$ (it is odd by
Lemma~\ref{lem:LB1}(iii), $x$ being odd in $\tilde G$), and since
$\deg_{\tilde G}(x)=\deg_G(x)-1\ge5$, the unique path ending at $x$ continues
through $x$; cutting it at $x$ adds one path (at most $\floor{n/2}$ in all)
and raises the endpoint count at $x$ to $3$ --- route~(b);
\item a decomposition $\tilde D$ with $|\tilde D|\le\floor{n/2}$ and $\tilde D(x)\ge3$: then
$xv$ is addible towards $v$, the passing count at $v$ being at most $1$ ($z$ is the only
neighbour of $v$ in $\tilde G$ that can be dead; all other neighbours are odd);
\item such a decomposition with $\tilde D(x)=1$ and $\tilde D(z)\ge1$: the passing count at
$v$ drops to $0$;
\item such a decomposition with $\tilde D(v)\ge2$ larger than the number of dead private
EV-neighbours of $x$: reinsert $xv$ towards $x$ (the passing count at $x$ sees only private
EV-neighbours of $x$, the other neighbours of $x$ being active), so that $v$ remains active
and the full star, including $uv$, is absorbable;
\item $\tilde G$ has at least two components of odd order: the component-wise
Theorem~\ref{thm:fan}, summed with Lemma~\ref{lem:sum}, gives route~(a).
\end{enumerate}
The budget arithmetic of the reinsertion skeleton behind routes (a)--(d),
and the formal statement that the exception of this corollary is exactly the
simultaneous failure of all five routes at \emph{every} common EV-neighbour,
are given in Appendix~\ref{app:band} (Theorem~\ref{thm:escape} and
Corollary~\ref{cor:corner}).
The stated exception is exactly the case in which, for every common EV-neighbour $v$, none of
the five routes applies.

\emph{Band closure of the corner.} One further narrowing is available in the constructed
family of Remark~\ref{rem:gap} (Definition~\ref{def:family}). Route~(a) can fail there only when $\tilde G$ carries a
\emph{deadlock band} (an even vertex of $\tilde G$ with more corridor neighbours than
terminal units; see Remark~\ref{rem:gap} for the definition), because a band-free $\tilde G$
in the family has $\pn(\tilde G)=\mathrm{odd}(\tilde G)/2$
(Theorem~\ref{thm:bandfree} in Appendix~\ref{app:band}); and the counting theorem
(Proposition~\ref{prop:gapnoband}) forces $\tilde G$ to be band-free whenever the counting gap
$\floor{n/2}-\mathrm{odd}(\tilde G)/2$ --- which is at least $1$ in this situation --- is at
most $2$, on the family proper (the constructed family without cross-private pairs).
Hence for every family-proper instance with counting gap at most $2$,
route~(a) applies unconditionally: the corner is empty at theorem level, not merely
computationally; the cross-private instances are backed at probe level by the machine
certificates of Appendix~\ref{app:comp}. For larger gaps the corner is blocked only by the excess bound recorded in
Remark~\ref{rem:gap}.
\end{proof}

\begin{corollary}\label{cor:uniquehigh}
If $G$ satisfies (A1) and (A2) and has \emph{exactly one} high vertex, then
$\pn(G)\le\floor{n/2}$.
\end{corollary}

\begin{proof}
Let $u$ be the unique high vertex. Every vertex other than $u$ is low, so every even vertex
$x\notin\{u\}\cup N_{\EV}(u)$ satisfies
$\dEV(x)-|N_{\EV}(x)\cap N_{\EV}(u)|\le\dEV(x)\le 3$; this is condition (A4) of
Theorem~\ref{thm:Astar}. (Alternatively, with only one high vertex the separation hypothesis of
Theorem~\ref{thm:sep} holds vacuously.)
\end{proof}

We now formulate the statement with neither the alignment nor the separation condition, which is
the form investigated in the computational study behind this paper and which all our computations
support.

\begin{conjecture}\label{conj:Astar}
Let $G$ be a connected graph such that $\EV(G)$ is triangle-free, and every vertex $v$ with
$\dEV(v)\ge 4$ has odd $\dEV(v)$ and only EV-neighbours $w$ with $\dEV(w)\le 3$. Then
$\pn(G)\le\floor{n/2}$.
\end{conjecture}

\begin{remark}[where the general argument stops]\label{rem:gap}
The class {\rm(A1)--(A2)} is not preserved by the only deletion the engine
supports. Deleting the E-star at a high vertex $u$ turns every star vertex odd;
an even vertex $x$ outside the star keeps its degree but loses the E-degrees
contributed by star vertices, and its E-degree in $\tilde G=G-F$ equals
$\dEV(x)-t_x$ with $t_x=|N_{\EV}(x)\cap N_{\EV}(u)|$. If $x$ is a high vertex
with $t_x$ odd and $\dEV(x)-t_x\ge 4$, then $x$ sits in $\EV(\tilde G)$ at an
even E-degree, the class is lost, and neither Theorem~\ref{thm:fan} nor the
induction applies to $\tilde G$. Concretely, in the double star on
$\{u,x,v_1,w_1,\dots,w_4,y_1,\dots,y_4\}$ with edges
$uv_1,uw_1,\dots,uw_4,xv_1,xy_1,\dots,xy_4$ and pendant leaves attaching so
that the eleven listed vertices have even degree, both $u$ and $x$ are high
with $\dEV=5$, $t_{ux}=1$ is odd, and after deleting the star at $u$ the vertex
$x$ has E-degree $4$, even. This graph (and its even-order one-leaf extension,
on which our exact programs give $\pn=5\le 11$) satisfies {\rm(A1)--(A2)} but
none of {\rm(A4)}, the separation condition, or the profile condition of
Remark~\ref{rem:CO}; a minimal counterexample to
Conjecture~\ref{conj:Astar}, if it exists, must have a high vertex with such a
bad partner at {\rm every} high vertex (Proposition~\ref{prop:residue}).
No counterexample is known: exact dynamic programs verify
$\pn\le\floor{n/2}$ on all tested instances satisfying {\rm(A1)--(A2)} (166
instances in the computational study behind this paper, a further $413$ in
Appendix~\ref{app:comp}, zero violations), and the family above sits far from
the bound ($\pn=\mathrm{odd}(G)/2$).

The ``one endpoint unit'' obstruction recorded in an earlier draft of this section was an
artefact: the reinsertion bookkeeping it refers to restores $k+2$ edges after deleting $k+1$,
i.e.\ counts the patched edge $xv_1$ twice (an activation step towards $v_1$ and a final step
towards $x$). Once $x$'s endpoint loss is recognised as cost-free --- $|D|$ is the only
invariant, and the argument of Theorem~\ref{thm:twohigh} involves no induction --- the
obstruction disappears for all two-high graphs except the corner recorded in
Corollary~\ref{cor:twohigh}.

The exceptional graphs of Corollary~\ref{cor:twohigh} are subject to five simultaneous
constraints, of which only the first two were recorded when the corner was first isolated:
the burial of the third EV-neighbour $z$ and of $v$ (or of enough private neighbours of $x$)
in \emph{every} optimal decomposition of $\tilde G$ (with $\tilde G$ as in
Corollary~\ref{cor:twohigh}), and the parity of the components of $\tilde G$, are each
individually sufficient to reopen the reinsertion argument (the edge $xv$ can then be added
towards $v$ at zero passing cost, or a second deletion $G-\tilde G$ becomes unnecessary). No
instance satisfying even the original two constraints was found: across $207$ constructed
instances of the dirty class---including the entire range where the counting gap
$\floor{n/2}-\mathrm{odd}(\tilde G)/2$ attains its minimum $1$, which the degree capacity of
the two far vertices forces to a single parameter-free shape---a decomposition of $\tilde G$
with at most $\floor{n/2}-1$ paths was always produced, so the corner never materializes in
computation.

The band-free tightness statement is now a theorem for the \emph{constructed family} (the
two-high graphs satisfying the invariants of our construction pipeline: no $\EV$-edge joins
two far even vertices of $\tilde G$; the odd-neighbour blocks of $u$ and of $x$ are of leaf
type (an odd number of pendant leaves) or gadget type, possibly mixed; every far vertex of
$\tilde G$ has degree at most $4$; the $207$ instances above and the $125$ decision-level
instances of Appendix~\ref{app:comp} all lie in it). Writing, for an even vertex $f$ of
$\tilde G$, $Y(f)$ for its \emph{corridor} neighbours --- the degree-two even vertices
adjacent to both $x$ and $f$ --- and $N(f)$ for its neighbour set in $\tilde G$, call $f$ a
\emph{deadlock band} if it has more corridor neighbours than terminal units, where the
terminal units of $f$ are its remaining incident edges (in symbols,
$|Y(f)|>|N(f)|-|Y(f)|$). If no even vertex is a deadlock band, then $\tilde G$ admits a
\emph{paired} decomposition: exactly $\mathrm{odd}(\tilde G)/2$ paths with all and only the
odd vertices of $\tilde G$ as endpoints. The proof (Theorem~\ref{thm:bandfree} in
Appendix~\ref{app:band}) is a five-step pairing
scheme---self-contained odd-neighbour blocks, forced fusions of the corridors and of the unit
at the common neighbour, an injective pairing of corridor ends with distinct terminal units
at every far vertex (which is precisely the negation of the band condition), and a sorted
offset matching at $x$ that never joins two corridors of the same far vertex, possible
because a far vertex of the family has degree at most $4$ and hence at most two corridors.
Every resulting trail visits $x$ once and leaves each far vertex for good, so no trail
closes. A single band, on the other hand, costs exactly one path: spending the free edge of
$x$ on a surplus corridor makes that corridor a double endpoint and lets its partner ride
through $x$ to a terminal, matching the exact value
$\pn(\tilde G)=\mathrm{odd}(\tilde G)/2+1$ computed by dynamic program. Consequently the
corner of Corollary~\ref{cor:twohigh} is empty whenever the counting gap is at most $2$,
unconditionally on the constructed family proper, with the cross-private extension covered
at probe level by the machine certificates of Appendix~\ref{app:comp} (Proposition~\ref{prop:gapnoband} and
Corollary~\ref{cor:bandiff} in Appendix~\ref{app:band}); for larger gaps it remains blocked only by the
excess bound $\mathrm{excess}\le\lfloor y/2\rfloor$, where
$\mathrm{excess}:=\pn(\tilde G)-\mathrm{odd}(\tilde G)/2$ and $y$ is the number of corridors
--- verified on all $235$ computed instances of the family, the provable bound being
$\pn(\tilde G)\le \mathrm{odd}(\tilde G)/2+2B-1$ for $B$ bands
(Proposition~\ref{prop:excb}), with a single band forcing
exactly $+1$ (Proposition~\ref{prop:patch}). The restriction to the constructed family is now
known to be sharper than necessary in two directions and irreducible in a third. Far--far
edges of the even subgraph and far vertices carrying several parity leaves, both excluded by
the construction, can be admitted: the corridor ends at a far vertex $f$ admit an injective
pairing with terminal units and far--far ends whenever $|Y(f)|\le t(f)+d(f)$ (where $t(f)$ is
the number of terminal units of $f$ and $d(f)$ its far--far degree), and the matching at $x$
survives arbitrary corridor multiplicities because a far vertex of the class has at most
three even neighbours, so a rearrangement argument replaces the sorted offset. What cannot be
admitted are arbitrary odd-neighbour blocks: a block whose component has exactly two odd
vertices admits a paired decomposition if and only if it admits a Hamilton path between them,
and the cycle-with-one-pendant and double-triangle blocks have none, so $\tilde G$ can be
untight while carrying no band at all---the band-free tightness statement is a phenomenon of
the constructed family, not of the class, and these cycle blocks are the first known
counterexamples to the general band-free statement. (Same-side pairs of private neighbours, a
third candidate obstruction, are excluded by the triangle-free condition {\rm(A1)} already.)
The counting gap of the dirty class also assumes that every even vertex lies in the
skeleton; blocks with even interior vertices shift $\floor{n/2}$ and must be recomputed
directly (see Appendix~\ref{app:comp} for the machine certificates).
\end{remark}

\begin{proposition}\label{prop:residue}
Let $G$ be a minimal counterexample to Conjecture~\ref{conj:Astar}. Then every
high vertex $u$ of $G$ admits a far high vertex $x$ with
$|N_{\EV}(x)\cap N_{\EV}(u)|$ odd and $\dEV(x)-|N_{\EV}(x)\cap N_{\EV}(u)|\ge4$.
In particular, by Theorem~\ref{thm:twohigh} a minimal counterexample has at least three high
vertices, unless it is one of the two-high exceptional graphs of Corollary~\ref{cor:twohigh}
(empty on the constructed family proper whenever the counting gap is at most two, by the
band-free tightness material of Remark~\ref{rem:gap} and Appendix~\ref{app:band}, with the
cross-private extension covered at probe level by the machine certificates of
Appendix~\ref{app:comp}; no instance was found in our computations, in or beyond the family).
\end{proposition}
\begin{proof}
Read ``minimal'' as minimal with respect to the number of edges, and let $u$
be a high vertex of $G$. We show that if $u$ has no bad partner, then
$\pn(G)\le\floor{n/2}$, contradicting the choice of $G$; hence $u$ has a bad
partner. Suppose, then, that no far even vertex $x$ satisfies
$t_x$ odd and $\dEV(x)-t_x\ge4$, where $t_x=|N_{\EV}(x)\cap N_{\EV}(u)|$.
Let $F$ be the E-star at $u$ and $\tilde G=G-F$. As in Step~1 of the proof of
Theorem~\ref{thm:Astar}, the even-degree vertices of $\tilde G$ are precisely
the far even vertices of $G$, $\EV(\tilde G)$ is an induced subgraph of
$\EV(G)$ and hence triangle-free, and a far vertex $x$ has E-degree
$\dEV(x)-t_x$ in $\tilde G$. Thus every high vertex of $\tilde G$ is a far
even vertex $x$ of $G$ with $\dEV(x)-t_x\ge4$; by our assumption $t_x$ is
even for such an $x$, and since $\dEV(x)\ge4$ makes $x$ high in $G$, (A2)
gives $\dEV(x)$ odd, so the E-degree $\dEV(x)-t_x$ of $x$ in $\tilde G$ is
odd. Moreover the EV-neighbours $w$ in $\tilde G$ of such an $x$ satisfy
$\dEV(w)-t_w\le\dEV(w)\le3$, the last inequality by (A2) applied to the high
vertex $x$ in $G$. Hence $\tilde G$ satisfies {\rm(A1)} and {\rm(A2)}. Each
component of $\tilde G$ is a connected graph satisfying {\rm(A1)}--{\rm(A2)}
with fewer edges than $G$ (as $u$ is high, $F\neq\emptyset$), so by minimality
each admits a decomposition with at most $\floor{n_c/2}$ paths, and
Lemma~\ref{lem:sum} gives a decomposition $\tilde D$ of $\tilde G$ with
$|\tilde D|\le\floor{n/2}$. Now the reinsertion argument of Steps~3 and~4 of
the proof of Theorem~\ref{thm:Astar} uses (A4) only through the existence of
such a $\tilde D$: the star vertices are odd in $\tilde G$, so the hypothesis
of Lemma~\ref{lem:fan7} holds read over $G$ exactly as in Step~3, and each
edge of $F\setminus A$ is addible towards its star vertex against a passing
count of at most $\dEV(v)-1\le2$ (by (A2)) and a budget of at least $3$. The
full star is restored with $|D|=|\tilde D|\le\floor{n/2}$ paths --- a
contradiction.
\end{proof}

We also record a structural consequence of (A1) on odd orders, used below to place the
sharpness family of Proposition~\ref{prop:oscsharp} relative to the class.

\begin{proposition}\label{prop:odddensity}
Let $G$ be connected of odd order $n$ with $\EV(G)$ triangle-free. Then
$m\le (n-1)^2/2$; equivalently, no odd semi-clique has triangle-free
$\EV(G)$.
\end{proposition}
\begin{proof}
Suppose $c=\binom{n}{2}-m\le (n-3)/2$. Let $S$ be the set of even-degree
vertices, $s=|S|$; since $n$ is odd, a vertex has odd degree iff its
complement-degree is odd, so each vertex outside $S$ has positive
complement-degree. Let $t$ be the number of complement edges inside $S$.
Since $G[S]=\EV(G)$ is triangle-free, Mantel's theorem gives
$\binom{s}{2}-t\le\floor{s^2/4}$, i.e.\ $t\ge \binom{s}{2}-\floor{s^2/4}$.
Counting complement degrees over $S$ and over its complement separately,
$2c\ge 2t+(n-s)$, so $2t\le s-3$. If $s$ is even this reads
$s(s-2)/2\le s-3$, i.e.\ $s^2-4s+6\le0$; if $s$ is odd,
$(s-1)^2/2\le s-3$, i.e.\ $s^2-4s+7\le0$; both are impossible. (Note
$s\ge1$: the number of odd-degree vertices is even while $n$ is odd.)
\end{proof}

We record the sharpness of the triangle-free requirement. By
Proposition~\ref{prop:odddensity} this failure is a density phenomenon: on odd orders, high
density and (A1) are mutually exclusive, so the following family is forced.

\begin{proposition}\label{prop:oscsharp}
For every $k\ge 2$, the graph $G_k=K_{2k+1}-(k-1)K_2$ (delete a matching of size $k-1$)
satisfies $\EV(G_k)\cong K_3$, hence violates (A1), and
$\pn(G_k)\ge k+1=\ceil{n/2}>\floor{n/2}$. Moreover $\pn(G_k)=k+1$ for the tested values
$k\in\{2,3,4,5,6,8,10\}$, with explicit decompositions machine-verified.
\end{proposition}

\begin{proof}
The $k-1$ matched vertices have degree $2k-1$ and the three unmatched vertices have degree
$2k$; the latter are pairwise adjacent, so $\EV(G_k)$ is a triangle on them. The edge count is
$m=\binom{2k+1}{2}-(k-1)=2k^2+1$, so Lemma~\ref{lem:LB2} gives
$\pn\ge\ceil{(2k^2+1)/2k}=k+1$; equivalently $G_k$ is an odd semi-clique and
Observation~\ref{obs:density} applies. This proves the necessity claim for all $k$. For the
tested values, explicit decompositions into $k+1$ paths (obtained from the Walecki cycle
decomposition of $K_{2k+1}$, spreading the matching edges over distinct cycles, cutting each
affected cycle at its matching edge and splitting the untouched cycle into a path plus a single
edge) were verified by machine to partition $E(G_k)$ exactly; for $k=2$ the value
$\pn(K_5-M_1)=3$ is also certified by exact dynamic programming.
\end{proof}

\begin{remark}
The exactness claim in Proposition~\ref{prop:oscsharp} is stated only for tested $k$: the
existence, for all $k$, of a Walecki-type cycle decomposition of $K_{2k+1}$ in which a given
$(k-1)$-matching is spread over distinct cycles is plausible but we do not prove it.
\end{remark}

Next, the witness family $W(m)$, which motivates Theorem~\ref{thm:Astar}: it lies in no
published class we are aware of, yet satisfies all hypotheses.

\begin{figure}[h]
\centering
\begin{minipage}{0.92\textwidth}\small
For odd $m\ge 5$, let $W(m)$ have vertex set
$\{a_1,a_2\}\cup\{b_1,\dots,b_m\}\cup\{\ell_1,\ell_2\}\cup\{h,h'\}\cup\{x_1,\dots,x_6\}$ and
edge set
\[
E(K_{2,m}) \;\cup\; \{a_1\ell_1,\,a_2\ell_2\} \;\cup\;
\{b_1h,\,b_1h'\} \;\cup\; \{hx_1,\,hx_2,\,h'x_1,\,h'x_2\} \;\cup\; E(K_6),
\]
where $K_6$ has vertex set $\{x_1,\dots,x_6\}$. Precisely: $h$ and $h'$ each have neighbourhood
$\{b_1,x_1,x_2\}$.
\end{minipage}
\caption{The witness family $W(m)$.}\label{fig:Wm}
\end{figure}

\begin{lemma}\label{lem:WmEV}
For every odd $m\ge 5$, the graph $W(m)$ of Figure~\ref{fig:Wm} satisfies
$\EV(W(m))\cong K_{2,m}$ exactly, hence satisfies (A1) and (A2) with $\dEV(a_1)=\dEV(a_2)=m$
and $\dEV(b_j)=2$ for all $j$.
\end{lemma}

\begin{proof}
The degrees are:
$\deg(a_i)=m+1$, even; $\deg(b_j)=2$ for $j\ge2$ and $\deg(b_1)=4$, all even;
$\deg(\ell_i)=1$, $\deg(h)=\deg(h')=3$, and $\deg(x_i)=5+2=7$ for $i\in\{1,2\}$,
$\deg(x_i)=5$ for $i\in\{3,\dots,6\}$, all odd. Hence the even-degree vertex set is exactly
$\{a_1,a_2\}\cup\{b_1,\dots,b_m\}$, and the induced subgraph is $K_{2,m}$: triangle-free, with
$\dEV(a_i)=m$ odd and all EV-neighbours of $a_i$ of E-degree $2\le3$. Note that the two
degree-$3$ vertices $h,h'$ attach to $b_1$ \emph{in pairs} (both $h$ and $h'$), keeping $b_1$
even; attaching a single such vertex would flip $b_1$ to odd degree and destroy the structure.
\end{proof}

\begin{proposition}\label{prop:Wmbound}
For every odd $m\ge 5$, $\pn(W(m))\le\floor{n/2}$ with $n=m+12$.
\end{proposition}

\begin{proof}
By Lemma~\ref{lem:WmEV} the graph satisfies (A1), (A2). We verify (A4) at $u=a_1$: the only
even-degree vertex outside $\{a_1\}\cup\{b_1,\dots,b_m\}$ is $a_2$, and, since
$N_{\EV}(a_1)=N_{\EV}(a_2)=\{b_1,\dots,b_m\}$,
$\dEV(a_2)-|N_{\EV}(a_2)\cap N_{\EV}(a_1)|=m-m=0\le3$; every vertex outside this set other than
$a_2$ has odd degree, hence is not in $\EV(\tilde G)$ after the deletion and is irrelevant (its
E-degree in $\tilde G$ is $0$). Thus (A4) holds, and
Theorem~\ref{thm:Astar} gives $\pn(W(m))\le\floor{(m+12)/2}$. For $m\in\{5,7,9\}$ explicit
decompositions into exactly $\floor{n/2}=8,9,10$ paths were in addition machine-verified.
\end{proof}

We now compare Theorem~\ref{thm:Astar} with the published record. The following two statements
make the incomparability precise; both witnesses have exactly computed path numbers.

\begin{proposition}[escape]\label{prop:escape}
Let $m\ge5$ be odd and let $W(m)$ be as in Figure~\ref{fig:Wm}. Then $\EV(W(m))\cong K_{2,m}$,
whence Theorem~\ref{thm:Astar} gives $\pn(W(m))\le\floor{n/2}$ with $n=m+12$. No published
theorem known to us implies this bound:
\begin{itemize}[leftmargin=1.5em]
\item every member of the Fan line \cite{Fan05,BS21,FHZ23} requires the blocks of $\EV(G)$ to
have maximum degree at most $3$, while $\Delta(\EV(W(m)))=m\ge5$;
\item $W(m)$ contains a $K_6$, so $\mathrm{tw}(W(m))\ge5$ and its degeneracy is at least $5$,
which excludes \cite{LY24} (treewidth $\le4$), \cite{BCLS20} (treewidth $\le3$) and
\cite{ZLH24} ($3$-degenerate);
\item $\Delta(W(m))=\max(m+1,7)\ge7$, which excludes \cite{BP19} ($\Delta\le5$) and
\cite{CFL} ($\Delta\le6$ with an independence condition);
\item $\Delta(W(m))\ge7$, so $W(m)$ is not subcubic and the polynomial-time algorithm of
\cite{MY26} for subcubic graphs does not apply;
\item $\EV(W(m))\cong K_{2,m}$ contains a $4$-cycle, so it is not a forest and \cite{Pyb96}
does not apply;
\item $W(m)$ is not Eulerian, so \cite{FK88} does not apply; it is not a block graph, so
\cite{CW25} does not apply.
\end{itemize}
For $m\in\{5,7,9\}$ every membership and non-membership claim above, together with explicit
decompositions into $\floor{n/2}$ paths, is machine-certified.
\end{proposition}

\begin{proposition}[proper inclusion in both directions]\label{prop:incomp}
The class of Theorem~\ref{thm:Astar} neither contains nor is contained in the Botler--Sambinelli
class $\mathcal{G}$ \cite{BS21} or the Fan--Hou--Zhou class \cite{FHZ23}.
\begin{enumerate}[label={\rm(\roman*)},leftmargin=2em]
\item Let $G$ be $K_{2,5}$ with one pendant leaf attached to each of the two degree-$5$
vertices. Then $\EV(G)\cong K_{2,5}$ is triangle-free, the two degree-$5$ vertices have odd
E-degree $5$ with all E-neighbours of E-degree $2$, so $G$ satisfies (A1)--(A2); but
$\Delta(\EV(G))=5$, so $G$ lies in none of \cite{Fan05,BS21,FHZ23}. Moreover
$\pn(G)=3\le\floor{9/2}$ by exact computation. (For the record, $G$ is $2$-degenerate, so its
bound is also implied by \cite{ZLH24}; this witness is used only for the incomparability with
the Fan line.)
\item Let $G'$ be a triangle $\{t_1,t_2,t_3\}$ together with two vertices $w_1,w_2$ adjacent
only to $t_1$ and to three pendant leaves each, so that $\EV(G')$ consists of the triangle
$\{t_1,t_2,t_3\}$ plus the two bridge edges $t_1w_1$, $t_1w_2$. Every block of $\EV(G')$ has
maximum degree at most $3$ and its single component contains exactly one block with triangles,
so $G'$ lies in $\mathcal{G}$ and in the Fan--Hou--Zhou class; yet the triangle in $\EV(G')$
violates (A1), so Theorem~\ref{thm:Astar} is silent on $G'$. Note $G'$ is not one of the
degenerate exceptions of the ``Gallai-or-SET'' alternative of \cite{BS21}. Exactly,
$\pn(G')=4\le\floor{11/2}$ by exact computation, while the main theorem of \cite{BS21}
gives only $\pn(G')\le\ceil{11/2}=6$ on its class (the Botler--Sambinelli bound is stated
with Gallai's ceiling, as noted in Section~\ref{sec:intro}); the Fan--Hou--Zhou membership
of $G'$ rests on the block structure of $\EV(G')$ alone and is unaffected.
\end{enumerate}
\end{proposition}

